\section{Method}
\label{sec:method}
\paragraph{Inference.}
We model standardised scores of a model as $x=\mu+\Lambda z+\varepsilon$ with $z\sim\mathcal N(0,I_r)$ and $\varepsilon\sim\mathcal N(0,\Psi)$, $\Psi=\mathrm{diag}(\psi_j)$ (factor analysis).
$\mu,\Lambda,\Psi$ are fitted by EM on the history, with missing cells handled by the marginal likelihood of the observed entries and a floor on $\psi_j$; the rank $r$ is chosen by held-out likelihood on history models.
Given the scores $x_S$ of the $k$ benchmarks run, the posterior of the latent profile is Gaussian with covariance $P_S=(I+\Lambda_S^\top\Psi_S^{-1}\Lambda_S)^{-1}$, the prediction for each target is $\hat x_j=\mu_j+\lambda_j^\top \mathbb{E}[z\mid x_S]$, and its predictive variance is $\psi_j+\lambda_j^\top P_S\lambda_j$.
The model is an assumption; we measure rather than assume its rank, and compare its accuracy with nonlinear imputers (Section~\ref{sec:experiments}).

\paragraph{Selection.}
Choosing $S$ is Bayesian experimental design on this model.
Because $P_S^{-1}=I+\sum_{j\in S}\lambda_j\lambda_j^\top/\psi_j$, adding benchmark $j$ to $S$ changes $P$ by a rank-one update, and the gain in the latent variance reduction $f_W(S)=\mathrm{tr}\,W(I-P_S)$ is in closed form,
\begin{equation}
\Delta_j(S)=\frac{\lambda_j^\top P_S W P_S\lambda_j}{\psi_j+\lambda_j^\top P_S\lambda_j},\qquad P_{S+j}=P_S-\frac{P_S\lambda_j\lambda_j^\top P_S}{\psi_j+\lambda_j^\top P_S\lambda_j},
\label{eq:gain}
\end{equation}
so a greedy step costs $O(Nr^2)$.
We use three weights: $W=\Lambda^\top\Lambda$ (all latent directions matter; \emph{A-optimal}); $W=\Lambda_T^\top\Lambda_T$ for a fixed target group $T$ (\emph{target-aware}, for transfer to a group of tasks not run); and the criterion that minimises the expected squared error on the complement of the selected set, in which the selected benchmarks leave the target set (\emph{complement-error}). 
A cost-aware variant (gain per unit of measured evaluation time) and a rank-free variant that conditions an empirical shrunk covariance directly complete the family.
The baselines are random choice, most-covered, centrality, k-medoids clustering of $|{\rm corr}|$, pivoted QR, leverage scores and greedy forward regression, all fitted on history rows.

\paragraph{Forecast and stopping.}
The posterior variance forecasts the error of a selection from the history fit alone, $\mathrm{MSE}(S)=\frac1{|T|}\bigl(\sum_{j\in T}\psi_j+\mathrm{tr}(\Lambda_T^\top\Lambda_T P_S)\bigr)$, with the finite-history correction of Section~\ref{sec:theory}.
Two stopping rules choose $k$: the \emph{model-based} rule takes the smallest $k$ at which the forecast RMSE is $\le\epsilon$; the \emph{split-conformal} rule \citep{lei2018distribution} takes the smallest $k$ at which the mean (or the 90\% quantile) of leave-one-model-out RMSE over history models is $\le\epsilon$, which needs no model assumption but assumes the new model resembles the history.

\paragraph{Applicability and cost.}
Computation is small next to running a benchmark; the cost that matters is the history, which must be large enough to resolve factors whose signal-to-noise is above $\sqrt{N/M}$ (Section~\ref{sec:theory}) and must contain the benchmark: one not in the history cannot be chosen or predicted.
